\section{Introduction}

Finite model clocks are useful only when their parameters can be traced back to the
physical arithmetic source.  The preceding prime-shell analysis used a clock of
modulus $4LQ$, with multiplier depth approximately $L$.  It produced exact support
theorems and a depth-uniform Bessel estimate after row normalization.  Those results
are internally correct, but two transfer questions remained open: whether one modeled
clock represents one physical denominator, and whether row normalization respects the
four-packet source identity.

This note resolves both questions at the structural level.  The physical row has two
linked scales, not one: its depth is $\lambda_h=hQ/H$, while its modulus is the
denominator $h$.  The earlier clock imposes a different modulus--depth relation.
Matching both relations is possible exactly at the exceptional scale $H=4Q^2$; V59
has $Q^2/H=x^{1/96}$ and therefore lies on the growing-mismatch side.

The normalization issue is independent.  A Bessel theorem for normalized vectors is
a legitimate frame-theoretic statement \citep{christensen2016}, but the physical
source is recovered by a signed four-phase polarization formula.  Normalizing each
output separately is nonlinear and makes that formula identically zero.  Any legal
conditioning transform must instead be common and linear across packet labels, with
all compensating weights retained explicitly.

The conclusions are deliberately scoped.  We prove an exact parameter crosswalk and
two incompatibility statements.  We do not prove cancellation in the divisor weights,
an $L^2$ saving, or the full Gate-B estimate.
